% !TEX program = lualatex
% Auto-generated from syllabus — 05: 知的財産と個人情報の基礎

\documentclass[aspectratio=43,professionalfonts,handout,10pt]{beamer}
\usepackage{beamer_template}

\newcommand{\LectureNum}{05}
\newcommand{\LectureTitle}{知的財産と個人情報の基礎}
\newcommand{\CourseName}{情報リテラシー}
\newcommand{\TextPages}{06-07}

\begin{document}
% --- Slide 1: title ---

\begin{frame}{\CourseName\quad 第\LectureNum 回「\LectureTitle 」}
  {\small \TermName\quad \TeacherName\quad ／\quad 教科書 \TextPages}

  \vfill

  \begin{block}{今日の流れ}
    \begin{enumerate}
      \item 知的財産権の種類と仕組み
      \item 著作権：何がOKで何がNGか
      \item 引用のルール
      \item 個人情報の定義と保護
    \end{enumerate}
  \end{block}
  \vfill

  \begin{block}{今日の目標}
    \begin{itemize}
      \item 知的財産権の種類（著作権・特許権・商標権等）を説明できる。
      \item 引用のルールと著作権上の許可・禁止事項を理解できる。
      \item 個人情報の定義と適切な取り扱いを説明できる。
    \end{itemize}
  \end{block}
  \vfill

  \begin{exampleblock}{キーワード}
    \small \textbf{著作権} ／ \textbf{産業財産権}（特許権・意匠権・商標権） ／ \textbf{引用} ／ \textbf{個人情報保護法} ／ \textbf{肖像権}
  \end{exampleblock}
\end{frame}

% --- Slide 2: 知的財産権の種類 ---

\begin{frame}{知的財産権の種類}
  \begin{description}[labelwidth=6em]
    \item[\kw{著作権}]
      文章・音楽・画像・プログラム等の創作物を保護する権利。創作と同時に自動的に発生する（\kw{無方式主義}）。保護期間は著作者の\kw{死後70年}
    \item[\kw{特許権}]
      新しい発明を保護する権利。出願・審査が必要（出願から\kw{20年}）
    \item[\kw{実用新案権}]
      物品の形状・構造などの考案を保護する権利（出願から\kw{10年}）
    \item[\kw{意匠権}]
      製品のデザイン（外観）を保護する権利（出願から\kw{25年}）
    \item[\kw{商標権}]
      商品・サービスのブランド名やロゴを保護する権利（登録から\kw{10年}，更新可）
  \end{description}

  \begin{alertblock}{演習1}
    著作権・特許権・商標権の違いを説明せよ
  \end{alertblock}

  \begin{slidenote}
    著作権は創作物（文章・音楽等）を保護し自動発生。特許権は新しい発明を保護し出願が必要。商標権はブランド名・ロゴを保護し更新が可能。保護対象・発生方法・期間の違いに注目。
  \end{slidenote}
\end{frame}

% --- Slide 3: 著作権 やっていいこと・いけないこと ---

\begin{frame}{著作権：やっていいこと・いけないこと}
  \begin{columns}[T]
    \begin{column}{0.48\textwidth}
      \begin{exampleblock}{許可されること（例）}
        \small
        \begin{itemize}
          \item 私的使用のための複製（自分だけが使う場合）
          \item 適切な引用（出典を明記し，引用部分が従）
          \item 教育機関での一定の利用（授業目的等）
        \end{itemize}
      \end{exampleblock}
    \end{column}
    \begin{column}{0.48\textwidth}
      \begin{alertblock}{禁止されること（例）}
        \small
        \begin{itemize}
          \item 他人の文章・画像を無断でWebに転載する
          \item 音楽・動画を無断でコピーして配布する
          \item 出典なしに他人の文章を自分のものとして使う（剽窃）
        \end{itemize}
      \end{alertblock}
    \end{column}
  \end{columns}

  \vfill
  \begin{slidenote}
    「私的使用」は自分だけの範囲。友人に送る・SNSに投稿する時点で私的使用の範囲外。教育目的の利用にも条件がある（授業で必要な範囲に限る）。
  \end{slidenote}
\end{frame}

% --- Slide 4: 引用のルール ---

\begin{frame}{引用のルール}
  \begin{block}{正しい引用の条件}
    \begin{itemize}
      \item 引用する\kw{必然性}がある（引用なしでは説明できない）
      \item 引用部分と自分の文章が\kw{明確に区別}されている
      \item 引用部分は\kw{従}，自分の文章が\kw{主}であること
      \item \kw{出典}（著者名・書籍名・URL・参照日）を明記する
    \end{itemize}
  \end{block}

  \vfill

  \begin{exampleblock}{レポート・スライドでの注意}
    Webの画像や文章を使う際は必ず出典を記載すること。「ネットで拾った画像」は著作物。
  \end{exampleblock}

  \begin{alertblock}{演習2}
    レポートで他者の文章を引用する際に守るべきルールを3つ挙げよ
  \end{alertblock}

  \begin{slidenote}
    引用の4条件を覚えること。特に「主従関係」が重要：自分の文章が主，引用は補足。引用ばかりのレポートはNG。Wikipediaも著作物なので出典記載が必要。
  \end{slidenote}
\end{frame}

% --- Slide 5: 個人情報の基礎 ---

\begin{frame}{個人情報の定義と保護}
  \begin{description}[labelwidth=8em]
    \item[\kw{個人情報}]
      生存する個人を特定できる情報。大きく2種類に分かれる
      \begin{itemize}
        \item 組み合わせで特定できる情報：氏名・住所・生年月日・電話番号・メールアドレス等
        \item 個人識別番号・身体的特徴：マイナンバー・免許証番号・顔や指紋のデータ等
      \end{itemize}
    \item[\kw{個人情報保護法}]
      個人情報の適切な取り扱いを定めた法律。企業・機関の義務を規定する
    \item[\kw{プライバシー権}]
      自分に関する情報を他人に知られたり勝手に公開されたりせず，\kw{自己決定できる}権利
    \item[\kw{肖像権}]
      自分の顔や姿を無断で撮影・公開されない権利
    \item[\kw{マイナンバー}]
      全住民に割り当てられた12桁の番号。特に厳格な管理が必要な\kw{特定個人情報}
  \end{description}

  \begin{slidenote}
    個人情報は「氏名」単独でも該当する。メールアドレスも個人を特定できれば個人情報。肖像権の侵害は法律上認められておらず，裁判でも保護が認められている。マイナンバーは特に厳格な管理が求められる「特定個人情報」に該当する。
  \end{slidenote}
\end{frame}

% --- Slide 6: 個人情報の取り扱い ---

\begin{frame}{個人情報の取り扱い上の注意}
  \begin{noticebox}[注意]
    \begin{itemize}
      \item SNSに氏名・住所・学校名・顔写真を公開する際は慎重に判断する
      \item 他人の個人情報を本人の同意なく第三者に伝えてはいけない
      \item パスワード・学籍番号などの認証情報は個人情報として厳重に管理する
      \item 事業者の管理不備による流出やフィッシングサイトなど，情報漏洩のリスクがある
    \end{itemize}
  \end{noticebox}

  \begin{alertblock}{演習3}
    SNSに写真を投稿する際に配慮すべき個人情報の観点を，具体例を挙げて述べよ
  \end{alertblock}

  \begin{slidenote}
    背景から住所特定，Exifの位置情報から撮影場所判明，他人の映り込みは肖像権侵害のおそれ，制服・校章から学校特定などのリスクがある。
  \end{slidenote}
\end{frame}

% --- チェックリスト ---

\begin{frame}{まとめ・チェックリスト}
  \begin{block}{自己チェック（できたら $\checkmark$ をつけよう）}
    \begin{itemize}
      \item[$\square$] 知的財産権の種類（著作権・特許権・意匠権・商標権）を説明できる
      \item[$\square$] 著作権で許可されること・禁止されることを理解した
      \item[$\square$] 引用のルール（4つの条件）を説明できる
      \item[$\square$] 個人情報の2分類と保護法の概要を説明できる
      \item[$\square$] プライバシー権・肖像権の意味を理解した
    \end{itemize}
  \end{block}

  \vfill

  {\small 質問があれば授業後またはTeamsで受け付けます。}
\end{frame}

\end{document}
